\documentclass[a4paper,12pt]{article}
\usepackage[T2A]{fontenc}
\usepackage[utf8]{inputenc}
\usepackage[russian]{babel}
\usepackage{amsmath,amssymb,mathtools}
\usepackage{PTSans}
\renewcommand{\familydefault}{\sfdefault}
\usepackage{icomma}
\usepackage[a4paper,margin=20mm,headsep=7mm,footskip=12mm]{geometry}
\usepackage{microtype}
\usepackage{tikz}
\usetikzlibrary{calc,arrows.meta,positioning,shapes.geometric}
\usepackage{tabularx,booktabs,longtable}
\usepackage{enumitem}
\usepackage{hyperref}
\usepackage{parskip}
\usepackage{fancyhdr}
\usepackage{setspace}
\usepackage{xcolor}

\definecolor{accent}{HTML}{2F6F73}
\definecolor{darkaccent}{HTML}{19484B}
\definecolor{textgray}{HTML}{2E3336}
\definecolor{softgray}{HTML}{6A7377}
\definecolor{rulegray}{HTML}{CBD4D6}
\hypersetup{colorlinks=true,linkcolor=darkaccent,urlcolor=darkaccent}
\color{textgray}
\onehalfspacing
\setlength{\parindent}{0pt}
\setlist[itemize]{leftmargin=6mm,itemsep=1.5mm,topsep=1.5mm}
\setlist[enumerate]{leftmargin=7mm,itemsep=2mm,topsep=2mm}
\setlength{\headheight}{15pt}
\setlength{\tabcolsep}{6pt}
\renewcommand{\arraystretch}{1.25}

\pagestyle{fancy}
\fancyhf{}
\fancyhead[L]{\small\color{softgray}Текстовые задачи на движение}
\fancyhead[R]{\small\color{softgray}04.10.26}
\fancyfoot[C]{\small\color{softgray}\thepage}
\renewcommand{\headrulewidth}{0pt}
\renewcommand{\footrulewidth}{0pt}

\newcommand{\sectiontitle}[1]{%
  \vspace{2mm}%
  {\Large\bfseries\color{darkaccent}#1}\par
  \vspace{1mm}\color{rulegray}\hrule\color{textgray}\vspace{3mm}%
}
\newcommand{\key}[1]{\textbf{\color{darkaccent}#1}}
\newcommand{\formula}[1]{\[\displaystyle #1\]}
\newcommand{\unit}[1]{\,\text{#1}}

\begin{document}

% ---------- TITLE ----------
\thispagestyle{empty}
\begin{center}
\vspace*{30mm}
{\small\color{softgray}ЧЕК-ЛИСТ}\par
\vspace{6mm}
{\Huge\bfseries\color{darkaccent}Текстовые задачи\\[2mm]на движение}\par
\vspace{8mm}
{\Large Табличный метод: от условия к уравнению}\par
\vspace{18mm}
{\large 04.10.26}\par
\vfill
{\large\bfseries Лёвин Артём Александрович}\par
\vspace{2mm}
{\large эксклюзивно для Матвея Горбачёва}\par
\vspace{18mm}
\end{center}

\begin{tikzpicture}[remember picture,overlay]
  \draw[accent,line width=1.3pt]
    ($(current page.south west)+(15mm,18mm)$)
    .. controls ($(current page.south west)+(65mm,32mm)$) and ($(current page.south east)+(-70mm,7mm)$) ..
    ($(current page.south east)+(-15mm,20mm)$);
  \draw[rulegray,line width=0.7pt]
    ($(current page.south west)+(15mm,13mm)$)
    .. controls ($(current page.south west)+(70mm,25mm)$) and ($(current page.south east)+(-65mm,2mm)$) ..
    ($(current page.south east)+(-15mm,15mm)$);
\end{tikzpicture}
\clearpage

% ---------- CORE ----------
\sectiontitle{1. Главная идея урока}

Большинство задач на движение удобно сводить к одной таблице:

\begin{center}
\begin{tabularx}{0.9\linewidth}{@{}lXXX@{}}
\toprule
Объект & Скорость $v$ & Время $t$ & Путь $S$ \\
\midrule
1 &  &  &  \\
2 &  &  &  \\
\bottomrule
\end{tabularx}
\end{center}

Связь трёх величин одна:
\formula{S=vt,\qquad v=\frac{S}{t},\qquad t=\frac{S}{v}.}

\key{Рабочий принцип.} Сначала внесите то, что дано напрямую. Затем восстановите третий столбец по формуле. После этого переведите словесную связь из условия в уравнение.

\key{Первый выбор переменной.} По возможности обозначайте через $x$ именно то, что требуется найти. Если этого недостаточно, вводите вспомогательную переменную: две переменные в таблице сами по себе не являются ошибкой.

\sectiontitle{2. Универсальный алгоритм}

\begin{enumerate}
  \item Определите объекты и тип движения: навстречу, вдогонку, по кругу, по воде, средняя скорость, протяжённые тела.
  \item Приведите единицы к системе, согласованной со скоростью. Если скорость в км/ч, время удобно выражать в часах, расстояние — в километрах.
  \item Постройте таблицу $v$--$t$--$S$.
  \item Внесите известные данные и выберите переменную.
  \item Заполните недостающий столбец формулами $S=vt$ или $t=S/v$.
  \item Найдите в тексте связь между полученными выражениями: равенство, сумма, разность, «на один круг больше», общее время и т. п.
  \item Составьте уравнение, решите его и вернитесь к вопросу задачи.
  \item Проверьте размерность и физический смысл ответа.
\end{enumerate}

\key{Особенно полезный ориентир:} если один объект стартовал раньше, к моменту встречи он двигался дольше. При сравнении времён или путей держите перед глазами правило «большее минус меньшее = известная разность».

\clearpage

% ---------- FLOWCHART ----------
\thispagestyle{plain}
\begin{center}
{\Large\bfseries\color{darkaccent}Алгоритм: текстовая задача на движение}
\end{center}
\vspace{6mm}

\begin{center}
\begin{tikzpicture}[
  node distance=7mm,
  every node/.style={font=\small,align=center},
  start/.style={ellipse,draw=accent,line width=1pt,minimum width=42mm,minimum height=10mm},
  action/.style={rectangle,rounded corners=2mm,draw=rulegray,line width=0.9pt,text width=105mm,minimum height=12mm,inner sep=3mm},
  branch/.style={rectangle,rounded corners=2mm,draw=rulegray,line width=0.9pt,text width=68mm,minimum height=16mm,inner sep=3mm},
  decision/.style={diamond,aspect=2.8,draw=accent,line width=1pt,text width=38mm,inner sep=1.5mm},
  arrow/.style={-{Stealth[length=2.3mm]},line width=0.9pt,color=softgray}
]
\node[start] (s) {Прочитать условие и вопрос};
\node[action,below=of s] (type) {Определить тип движения и привести единицы};
\node[action,below=of type] (table) {Составить таблицу $v$--$t$--$S$ и внести известные данные};
\node[action,below=of table] (var) {Выбрать $x$; при необходимости ввести вспомогательную переменную};
\node[action,below=of var] (formula) {Заполнить недостающий столбец через $S=vt$ или $t=S/v$};
\node[decision,below=9mm of formula] (link) {Какая связь дана в тексте?};
\node[branch,below=13mm of link,xshift=-42mm] (eq1) {Равные величины / общее время / одинаковый путь};
\node[branch,below=13mm of link,xshift=42mm] (eq2) {Разность / отрыв / один круг / задержка старта};
\node[action,below=40mm of link] (solve) {Составить и решить уравнение};
\node[action,below=of solve] (check) {Проверить единицы, знак, смысл и ответить именно на вопрос};
\node[start,below=of check] (e) {Ответ};

\draw[arrow] (s)--(type);
\draw[arrow] (type)--(table);
\draw[arrow] (table)--(var);
\draw[arrow] (var)--(formula);
\draw[arrow] (formula)--(link);
\draw[arrow] (link)--(eq1);
\draw[arrow] (link)--(eq2);
\draw[arrow] (eq1.south) |- (solve.west);
\draw[arrow] (eq2.south) |- (solve.east);
\draw[arrow] (solve)--(check);
\draw[arrow] (check)--(e);
\end{tikzpicture}
\end{center}
\clearpage

% ---------- TYPES ----------
\sectiontitle{3. Встречное движение и разновременный старт}

Если автомобили движутся навстречу и второй выехал позже на $\Delta t$, то первый к моменту встречи находится в движении на $\Delta t$ дольше.

\textbf{Пример.} Между городами 435 км. Первый автомобиль едет со скоростью $60$ км/ч, второй — $65$ км/ч и выезжает на 1 час позже. Пусть $x$ — путь первого автомобиля до встречи. Тогда
\formula{\frac{x}{60}-\frac{435-x}{65}=1.}
Здесь не нужно запоминать специальную формулу: уравнение непосредственно выражает разницу времён.

\sectiontitle{4. Движение вдогонку}

При движении в одном направлении быстрый объект сокращает отрыв со скоростью
\formula{v_{\text{сбл}}=v_{\text{быстр}}-v_{\text{медл}}.}
Если начальный отрыв равен $d$, то
\formula{t=\frac{d}{v_{\text{быстр}}-v_{\text{медл}}}.}

\textbf{Показательный случай.} Отрыв 300 м = $0{,}3$ км, разность скоростей $1{,}5$ км/ч:
\formula{t=\frac{0{,}3}{1{,}5}=0{,}2\text{ ч}=12\text{ мин}.}

\key{Контроль единиц.} Если скорости заданы в км/ч, полученное время сначала выражено в часах. Для перевода в минуты умножайте на 60.

\sectiontitle{5. Движение по круговой трассе}

Если два объекта стартовали \textbf{из одной точки} и движутся в одном направлении, то обгон на один круг означает:
\formula{S_{\text{быстр}}-S_{\text{медл}}=L,}
где $L$ — длина круга.

\textbf{Пример.} Длина круга 14 км, скорость быстрого автомобиля $80$ км/ч, время $40$ минут $=\frac23$ часа. Если скорость второго $x$, то
\formula{80\cdot\frac23-x\cdot\frac23=14,}
откуда $x=59$ км/ч.

\clearpage

\sectiontitle{6. Движение по воде}

Пусть $v$ — собственная скорость судна, $u$ — скорость течения. Тогда
\formula{v_{\downarrow}=v+u,\qquad v_{\uparrow}=v-u.}
Для плота собственная скорость равна нулю, поэтому его скорость относительно берега равна скорости течения.

Если судно прошло одинаковое расстояние $x$ по течению и против течения, а также стояло $t_0$, то общее время:
\formula{\frac{x}{v+u}+t_0+\frac{x}{v-u}=T.}

\textbf{Пример.} $v=25$ км/ч, $u=3$ км/ч, стоянка 5 часов, весь рейс занял 30 часов:
\formula{\frac{x}{28}+5+\frac{x}{22}=30.}
Получается $x=308$ км в одну сторону, полный путь — $616$ км.

\sectiontitle{7. Средняя скорость}

Всегда используйте определение:
\formula{v_{\text{ср}}=\frac{S_{\text{общ}}}{t_{\text{общ}}}.}

\textbf{Если равны доли пути,} среднюю скорость нельзя находить обычным средним арифметическим скоростей.

Например, треть пути пройдена со скоростью 12 км/ч, треть — 16 км/ч, треть — 24 км/ч. При общем пути $x$:
\formula{t_{\text{общ}}=\frac{x}{36}+\frac{x}{48}+\frac{x}{72}=\frac{x}{16},\qquad v_{\text{ср}}=16\text{ км/ч}.}

\textbf{Если равны доли времени,} средняя скорость действительно равна среднему арифметическому соответствующих скоростей.

\clearpage
\sectiontitle{8. Протяжённые тела}

Поезд, баржа, сухогруз и подобные объекты нельзя всегда считать материальными точками: для полного прохождения одного объекта мимо другого нужно учесть их длины.

Пусть длины объектов $L_1$ и $L_2$, а начальный промежуток между ними — $d$.

\textbf{В одном направлении:}
\formula{t=\frac{d+L_1+L_2}{v_{\text{быстр}}-v_{\text{медл}}}.}

\textbf{Навстречу:}
\formula{t=\frac{d+L_1+L_2}{v_1+v_2}.}

Если начального промежутка нет, берите $d=0$.

\key{Важно.} Сначала приведите длины и скорости к согласованным единицам: например, метры и м/с либо километры и км/ч.

% ---------- ERRORS ----------
\sectiontitle{9. Типичные ошибки}

\begin{enumerate}[itemsep=1mm]
  \item \textbf{Неверный порядок разности.} Кто стартовал раньше, тот к моменту встречи двигался дольше; при обгоне быстрый проходит больше.
  \item \textbf{Смешение единиц.} Метры переводите в километры при скоростях в км/ч; минуты — в часы. Например, 300 м $=0{,}3$ км, 40 минут $=\frac23$ часа.
  \item \textbf{Потеря условия типа движения.} По воде используйте $v\pm u$; на круге учитывайте длину круга и точку старта.
  \item \textbf{Неверная средняя скорость.} Основная формула — $v_{\text{ср}}=S_{\text{общ}}/t_{\text{общ}}$; среднее арифметическое подходит только в специальных случаях, например при равных промежутках времени.
  \item \textbf{Протяжённые тела приняли за точки.} Для полного обгона или расхождения учитывайте сумму длин и, если он есть, начальный промежуток.
  \item \textbf{Решили лишнее или ответили не на вопрос.} Ищите требуемую величину как можно прямее и в конце проверяйте размерность и смысл.
\end{enumerate}

\sectiontitle{10. Быстрая самопроверка перед ответом}

\textbf{Перед ответом проверьте:} единицы согласованы; таблица подписана; уравнение точно переводит ключевую фразу условия; специальные особенности типа движения учтены; корень физически допустим; финальная запись отвечает именно на вопрос и содержит нужные единицы.

\clearpage

% ---------- TRAINING ----------
\sectiontitle{11. Тренировочный блок — Путь А}

\textbf{10 задач.} Решайте табличным методом. В каждом задании сначала укажите единицы, затем составьте таблицу и уравнение. Решения в этот блок не записаны.

\begin{enumerate}
  \item Расстояние между городами $A$ и $B$ равно 340 км. Из $A$ в $B$ выехал автомобиль со скоростью 60 км/ч. Через 1 час навстречу ему из $B$ выехал второй автомобиль со скоростью 80 км/ч. На каком расстоянии от $A$ они встретятся?

  \item Медленный автомобиль движется со скоростью 54 км/ч. Быстрый автомобиль, находясь позади него на 9 км, движется в том же направлении со скоростью 72 км/ч. Через сколько минут быстрый автомобиль догонит медленный?

  \item Грузовик выехал из города со скоростью 50 км/ч. Через 2 часа вслед за ним выехал легковой автомобиль. Он догнал грузовик через 4 часа после своего старта. Найдите скорость легкового автомобиля.

  \item Длина круговой трассы 15 км. Из одной точки одновременно в одном направлении стартовали два автомобиля. Скорость быстрого автомобиля 90 км/ч. Через 20 минут он опередил второй автомобиль ровно на один круг. Найдите скорость второго автомобиля.

  \item Собственная скорость лодки 20 км/ч, скорость течения 4 км/ч. Лодка прошла некоторое расстояние по течению, затем стояла 1 час и вернулась тем же путём против течения. Весь рейс занял 6 часов. Найдите полный путь лодки.

  \item Плот движется по течению со скоростью 3 км/ч. Через 2 часа после его старта из той же точки вслед за ним отправилась моторная лодка с собственной скоростью 15 км/ч. Через сколько минут после своего старта лодка догонит плот?

  \item Первую треть пути автомобиль проехал со скоростью 12 км/ч, вторую треть — 18 км/ч, последнюю треть — 36 км/ч. Найдите среднюю скорость на всём пути.

  \item Первую треть времени автомобиль ехал со скоростью 30 км/ч, вторую треть времени — 60 км/ч, последнюю треть времени — 90 км/ч. Найдите среднюю скорость за всё время движения.

  \item Два поезда длиной 180 м и 120 м движутся в одном направлении. Скорость более быстрого поезда 72 км/ч. Он полностью обгоняет второй поезд за 30 секунд. Найдите скорость более медленного поезда.

  \item Два поезда длиной 150 м и 250 м движутся навстречу друг другу со скоростями 54 км/ч и 90 км/ч. Сколько секунд проходит от момента встречи их головных вагонов до полного расхождения поездов?
\end{enumerate}

\clearpage

% ---------- ANSWERS ----------
\sectiontitle{12. Ответы}

\begin{center}
\begin{tabularx}{0.92\linewidth}{@{}>{\bfseries}cX@{}}
\toprule
№ & Ответ \\
\midrule
1 & $180$ км от города $A$ \\
2 & $30$ минут \\
3 & $75$ км/ч \\
4 & $45$ км/ч \\
5 & $96$ км \\
6 & $24$ минуты \\
7 & $18$ км/ч \\
8 & $60$ км/ч \\
9 & $36$ км/ч \\
10 & $10$ секунд \\
\bottomrule
\end{tabularx}
\end{center}

\vspace{7mm}
\textbf{Ключевые проверки:}
\begin{itemize}
  \item №1: $\frac{x}{60}-\frac{340-x}{80}=1$.
  \item №4: $(90-v)\cdot\frac13=15$.
  \item №5: $\frac{x}{24}+1+\frac{x}{16}=6$, где $x$ — путь в одну сторону; в ответ требуется $2x$.
  \item №7: при общем пути $S$ время равно $\frac{S}{36}+\frac{S}{54}+\frac{S}{108}=\frac{S}{18}$.
  \item №9: относительная скорость $\frac{180+120}{30}=10$ м/с $=36$ км/ч.
  \item №10: относительная скорость $54+90=144$ км/ч $=40$ м/с; общая длина $400$ м.
\end{itemize}

\vfill
{\small\color{softgray}Лёвин Артём Александрович эксклюзивно для Матвея Горбачёва}

\end{document}
